\section{Relation with Actual Causality} \label{sec:relation_with_ac}

The notion of explanation-extended distribution is not meant to explicate Halpern's notion of actual
causality; rather, the goal was to combine the intuitions underlying Halpern's notion with those
motivating Pearl's notion of probability of necessity and sufficiency to provide a more general
tool for causal attribution. Nevertheless, to show that we preserve the spirit of Halpern's
definition, in this section, we discuss the connection between actual causality and the notions of
explanation-extended distribution (for proofs, see the appendix). The goal is to show that there is
a natural choice of parameters for our method that makes it give verdicts that agree with actual
causality. With one caveat: probability of causation maximization will not identify subset-minimal
causes that are not cardinality-minimal.



Let's start with the main definition used by Halpern:

\begin{definition} \label{def:ac}
    \(\mathbf{C} = \mathbf{c}\) is an actual cause of \(\varphi\) in the causal setting \((M, \mathbf{u})\) if the following three conditions hold:
    \begin{itemize}
        \item \textbf{Factivity}: \((M, \mathbf{u}) \vDash (\mathbf{C} = \mathbf{c}\,)\) and \((M, \mathbf{u}) \vDash \varphi\),
        \item \textbf{Context-sensitive Necessity}: There is a set of variables \(\mathbf{T} \subseteq \mathbf{V}\) and an alternative setting \(\mathbf{c'}\) of variables in \(\mathbf{C}\) such that if \((M, \mathbf{u}) \vDash \mathbf{T} = \mathbf{t^\star}\), then
            \[(M, \mathbf{u}) \vDash [\mathbf{C} \leftarrow \mathbf{c'}, \mathbf{T} \leftarrow \mathbf{t^\star}] \neg \varphi,\]
        \item \textbf{Minimality}: \(\mathbf{C}\) is minimal; there is no strict subset \(\mathbf{C_{sub}} \subset \mathbf{C}\) such that \(\mathbf{C_{sub}} = \mathbf{c_{sub}}\) satisfies the above two conditions where $\mathbf{c_{sub}}$ is the restriction of \(\mathbf{c}\) to the variables in \(\mathbf{C_{sub}}\).
    \end{itemize}
\end{definition}


First, we notice that if the actual causality holds, there are some fairly natural settings of
suspects and witnesses such that for any (regular) witness and suspect biases, the resulting
integrated probability of causality (still parameterized by noise) is greater than zero.


\begin{theorem}\label{th:ac-exp}
    Suppose that $\mathbf{C}$ is an actual cause with a witness set $\mathbf{T}$ in $(M, \mathbf{u})$.  Take any  $\{\mathbf{S} = \mathbf{s}\} \supseteq \mathbf{C}$ as the potential cause set, and follow the definition in taking the set of all endogenous variables $\mathbf{V} \supseteq \mathbf{T}$ as the potential witness set. Then for any suspect bias  $ - 0.5 < b_s < 0.5$, and  any witness  bias $ - 0.5 < b_w < 0.5$ we have: \begin{align}
    P^{E}_{b_s, b_w , \varphi, \{\mathbf{S} = \mathbf{s}\}, \mathbf{V}}(\mathbf{u}, \mathbf{v}, C, T) > 0 \label{eq:pect}\\
    \text{ and therefore } \nonumber\\
\sum_{\mathbf{T} \subseteq \mathbf{V}}\, P^E_{b_s, b_w , \varphi, \{\mathbf{S} = \mathbf{s}\}, \mathbf{V}}(\mathbf{u}, \mathbf{v}, C, T) > 0 \end{align}
\end{theorem}


There is also a clear connection in the opposite direction: if the same noise-parameterized
probability of causation is greater than zero, the necessity condition in the definition of actual
causality holds.


\begin{theorem}\label{th:exp-ac} Moreover, if
$\sum_{T \subseteq \mathbf{V}} P^{E}_{b_s, b_w, \varphi, \mathbf{C}, \mathbf{V}}(\mathbf{u}, \mathbf{v}, \mathbf{C}, \mathbf{T}) > 0$,  the context-sensitive necessity condition in Definition \ref{def:ac} holds. \label{th:if_pr_then_nec}
\end{theorem}






With minimality, the situation is somewhat more complicated. For one thing, our choice of suspect
and witness preemption biases will have an impact on probability-preferred set sizes. On the other
hand, the way that  $P^{\mathit{set}}_{b_s,b_w, \{\mathbf{S}=\mathbf{s}\}, \mathbf{V}}(\mathbf{C}, \mathbf{T})$ is defined makes sets comparable in light of their cardinalities, whereas the minimality Definition \ref{def:ac}  is formulated in terms of the subset relation, which is bound to produce some misalignment. Some useful observations can be made, though.

One of them is that there is a somewhat degenerate simplified setting in which positive
explanation-extended probability, on the assumption of factivity, entails minimality wrt. to the
same parameters: it is enough to assume null witness bias and that the suspect bias prefers smaller
explanations. The strategy used in the proof entails that we can safely marginalize out the
witnesses as well.


\begin{theorem} \label{th:local_max_to_ac} Write $\Phi (\mathbf{C}, \mathbf{T})$ for $P^{E}_{b_s, b_w, \varphi, \mathbf{C}, \mathbf{V}}(\mathbf{u}, \mathbf{v}, \mathbf{C}, \mathbf{T})$. Consider the set of optimizers  $\Gamma = \arg\max_{\mathbf{C} \subseteq \{\mathbf{S} = \mathbf{s}\}, \mathbf{T} \subseteq \mathbf{V}} \Phi(\mathbf{C}, \mathbf{T})$.  Assume $(\mathbf{C^\dagger}, \mathbf{T^\dagger})\in \Gamma$ with $\Phi(\mathbf{C^\dagger}, \mathbf{T^\dagger})>0$ and assume further that factivity is satisfied for $\mathbf{C^\dagger}$ (as in Definition \ref{def:ac}). Let the witness bias be $b_w=0$ and let the suspect preemption bias be $b_s>0$ (that is, we positively prefer suspect preemption, thus building preference for smaller cause sets into our search). Then $C^\dagger$ is an actual cause of $\varphi$ according to Definition \ref{def:ac}.
\end{theorem}



\begin{corollary}\label{cor:marginalize} Let's keep the notation from Theorem \ref{th:local_max_to_ac}. Consider the set   of optimizers  $\Delta = \arg\max_{\mathbf{C} \subseteq \{\mathbf{S} = \mathbf{s}\}} \sum_{\mathbf{T} \subseteq \mathbf{V}} \Phi(\mathbf{C}, \mathbf{T})$.  Assume $\mathbf{C^\dagger}\in \Delta$ with \[\sum_{\mathbf{T} \subseteq \mathbf{V}} \Phi(\mathbf{C^\dagger}, \mathbf{T})>0\] and assume still that factivity is satisfied for $\mathbf{C^\dagger}$. Keep $b_w=0$ and $p_s>0$  as in Theorem \ref{th:local_max_to_ac}. Then $\mathbf{C^\dagger}$ is an actual cause of $\varphi$ according to Definition \ref{def:ac}.
\end{corollary}



Now we can return to the misalignment we have mentioned and elaborate on the nuance. One can
imagine a reasoning mode when one looks for causes of an outcome by inspecting probabilities of
causation $P(C \rightarrow \varphi)$ not knowing what the values of $\mathbf{u}$ are, but knowing that $\varphi$ occurred. It is worth pointing out that  $\arg\max_{\mathbf{C} \subseteq \mathbf{X}}P(\mathbf{C} \rightarrow \varphi)$ might fail to include some $\mathbf{C}$ such that for some setting $\mathbf{u}$, $\mathbf{C}$ is nevertheless an actual cause of $\varphi$.


\begin{observation}\label{obs:misalignment} Define $\Xi = \arg\max_{\mathbf{C} \subseteq {X}} P(C \rightarrow \varphi)$. It is not the case that if $C^\dagger$ is an actual cause in $(M, \mathbf{u})$ for some $\mathbf{u}$, $C^\dagger \in \Xi$.
\end{observation}

The rather straightforward reason for this is that actual cause sets might be subset-minimal
without being cardinality-minimal, while the probability of causality maximization will drop those
that are not cardinality-minimal.

This can be relatively easily alleviated by not maximizing the probability of causation if one
searches for actual causes, but rather finding all cause sets with non-zero probabilities (by our
previous theorems, all of those should satisfy the necessity requirement) and subsequently
filtering for subset-minimality.

However, we find that some idiosyncratic components of the notion of actual causality don't make
it directly scalable anyway, so at this point, we wrap up a discussion of that notion and move to
the discussion of a related, more probabilistic, and more scalable notion.


% \begin{observation} Consider
% \begin{align*}\Xi & = \left\{C \vert \left(\sum_{\mathbf{T} \subseteq \mathbf{V}} \Phi(C, T)>0\right),  \neg \exists C' \subset C \left( \sum_{\mathbf{T} \subseteq \mathbf{V}} \Phi(C', T)\right) \right\}
% \end{align*}

% \begin{enumerate}
% \item $\Delta \subseteq \Xi$.
% \item If $C^\dagger \in \Xi$, $C^\dagger$ is an actual cause.
% \end{enumerate}
% \end{observation}



% \begin{proof}
% TODO
% \end{proof}


% \subsection{Factivity}

% Our definition of explanation-extended distribution does not enforce the criteria of factivity but that can be easily achieved by conditioning. We also believe that the factivity is a restricting condition since we do not wish to limit our causal reasoning to a particular factual world. We wish to generalize to other possible worlds as well.



% \section{Questions}
% \begin{itemize}
%     \item These definitions do not exactly coincide with Halpern and Pearl's definitions. How do we want to go about justifying these definitions?
%     \item Are there any other quantities that people think we can define above?
% \end{itemize}
